\section{Vectorisation}

\subsection{Stratégie de Vectorisaton}

\subsubsection{Motivation}

Le point 5.1 de l'énoncé définit la \emph{vectorisation} comme un changement de représentation des agents :
au lieu de manipuler une liste d'objets « fourmi », on stocke séparément les attributs (coordonnées, état,
graine de générateur pseudo-aléatoire, etc.) dans des tableaux. Une fourmi est alors repérée uniquement par
son indice $i$ dans ces tableaux.

Cette étape est motivée par l'analyse \emph{baseline} : le bloc $K1$ (avance de l'ensemble des fourmis)
concentre l'essentiel du temps noyau. L'objectif principal est donc d'améliorer l'efficacité mémoire de cette
boucle (localité, cache, bande passante) afin de réduire le coût par itération, tout en préparant la suite
(parallélisation OpenMP au point 5.2 de la consigne).

\subsubsection{Du layout AoS au layout SoA}

Dans la version \emph{baseline}, la population est stockée sous la forme d'un \texttt{std::vector<ant>}
(\textit{Array of Structs}, AoS). Chaque instance \texttt{ant} regroupe plusieurs champs hétérogènes
(position \texttt{SDL\_Point}, état, graine RNG, etc.). Lors de l'itération sur toutes les fourmis, le noyau
lit/écrit principalement la position, l'état et la graine, mais les lignes de cache chargées contiennent aussi
des données non nécessaires au pas courant.

Nous introduisons donc un layout \textit{Structure of Arrays} (SoA), plus favorable à la localité et à la
vectorisation au sens « données » du sujet :
\[
\text{AoS: } \{(x_i, y_i, s_i, g_i)\}_{i=0}^{N-1}
\qquad\longrightarrow\qquad
\text{SoA: } (x[\,], y[\,], s[\,], g[\,]).
\]
Ainsi, les accès séquentiels \texttt{x[i]}, \texttt{y[i]}, \texttt{state[i]}, \texttt{seed[i]} deviennent
contigus en mémoire, ce qui augmente la proportion de données utiles par ligne de cache et facilite les boucles
par indice.

\subsubsection{Maintien du baseline : un binaire, deux layouts}

Afin de conserver une référence expérimentale stable, le chemin AoS est maintenu et reste le mode d'exécution
par défaut. Le programme expose un paramètre de ligne de commande permettant de choisir la représentation :
\begin{itemize}[leftmargin=*]
    \item \texttt{--layout aos} : exécution \emph{baseline} (\texttt{std::vector<ant>}) ;
    \item \texttt{--layout soa} : exécution vectorisée (\texttt{AntsSoA}).
\end{itemize}
Ce choix « 1 binaire, 2 layouts » évite de dupliquer le projet en deux répertoires (et donc de diverger au
fil des modifications) et garantit une comparaison directe AoS vs SoA à paramètres identiques.

\subsubsection{Choix d'implémentation : types compacts et seed par fourmi}

En SoA, nous stockons explicitement les attributs minimaux requis par la dynamique :
\begin{itemize}[leftmargin=*]
    \item \texttt{x[i]}, \texttt{y[i]} en \texttt{int32\_t} (coordonnées) ;
    \item \texttt{state[i]} en \texttt{uint8\_t} (0 = \textit{unloaded}, 1 = \textit{loaded}) ;
    \item \texttt{seed[i]} en \texttt{uint32\_t} (état RNG local).
\end{itemize}
Le choix de types compacts réduit l'empreinte mémoire par agent et améliore l'utilisation du cache, ce qui est
particulièrement pertinent puisque la boucle $K1$ parcourt plusieurs milliers de fourmis.

Un point important pour la reproductibilité est la gestion de la graine RNG par fourmi. Dans le code AoS
\emph{baseline}, le champ \texttt{m\_seed} est utilisé dans \texttt{ant::advance} (\texttt{src/ant.cpp}) mais
n'était pas initialisé dans le constructeur (\texttt{src/ant.hpp}), ce qui induit un \emph{comportement indéfini}
et rend le baseline instable. La vectorisation impose de rendre cette donnée explicite : chaque fourmi possède
sa graine \texttt{seed[i]} et les fonctions de génération aléatoire la mettent à jour par référence. Pour garder
une comparabilité maximale, la formule du générateur est conservée ; seule l'initialisation devient correcte et
déterministe.

\subsubsection{Application dans le noyau : \texttt{advance\_time\_soa}}

La version SoA introduit un noyau \texttt{advance\_time\_soa} (analogue à \texttt{advance\_time} dans
\texttt{src/ant\_simu.cpp}) :
\begin{enumerate}[leftmargin=*]
    \item boucle sur $i \in [0, N)$ : avance de la fourmi $i$ en manipulant \texttt{x[i]}, \texttt{y[i]},
    \texttt{state[i]}, \texttt{seed[i]} ;
    \item évaporation globale des phéromones (\texttt{do\_evaporation}) ;
    \item validation/swap de la carte de phéromones (\texttt{update}).
\end{enumerate}
La logique par fourmi reste volontairement identique à la version AoS (mêmes choix de voisin, mêmes branches,
même condition de mouvement \texttt{consumed\_time < 1}) afin de comparer les performances en isolant l'effet du
layout mémoire.

\subsubsection{Découplage de \texttt{SDL\_Point} hors hot loop et adaptation du rendu}

Le type \texttt{SDL\_Point} est pratique pour le rendu mais il couple le noyau de calcul à SDL. Dans la stratégie
SoA, le noyau manipule des scalaires \texttt{int32\_t} ; la conversion vers un point SDL (si nécessaire) est
repoussée au moment du rendu. Dans le même esprit, une surcharge \texttt{mark\_pheronome(x,y)} côté
\texttt{pheronome} permet de marquer les phéromones sans construire un \texttt{SDL\_Point} à chaque pas.

Enfin, le \texttt{Renderer} est adapté pour accepter soit une population AoS (\texttt{std::vector<ant>}), soit une
population SoA (\texttt{AntsSoA}). Le rendu itère alors :
\begin{itemize}[leftmargin=*]
    \item en AoS : sur les objets et \texttt{ant.get\_position()} ;
    \item en SoA : sur l'indice $i$ et les tableaux \texttt{x[i]/y[i]}.
\end{itemize}

\subsubsection{Validation et protocole de comparaison}

La validation repose sur la compilation en mode \emph{release} puis l'exécution du mode \texttt{--benchmark} en
forçant successivement \texttt{--layout aos} puis \texttt{--layout soa}. Le format des métriques
\texttt{METRIC ...} est conservé afin de réutiliser le pipeline de consolidation existant
(\texttt{scripts/measure\_temps.sh}). La comparaison se focalise sur $K1$ (et donc $K0$), puisque la vectorisation
vise prioritairement le noyau d'avance des fourmis.
